In its general form, the Fokker-planck equation may be expressed in
the conservative form\[
\frac{\partial f_{0}}{\partial t}+\nabla\cdot\mathbf{S}\left(f_{0}\right)=0\]
 where the flux $\mathbf{S}$ is\[
\mathbf{S}=-\overline{\overline{D}}\nabla f_{0}+\overline{\overline{F}}f_{0}\]
 $\overline{\overline{D}}$ and $\overline{\overline{F}}$ being respectively
the diffusion and convection tensors. One may express $\overline{\overline{D}}$
and $\overline{\overline{F}}$ in the general matrix form \[
\overline{\overline{D}}=\left(\begin{array}{llll}
D_{pp} & D_{p\xi} & D_{pr} & D_{p\theta}\\
D_{\xi p} & D_{\xi\xi} & D_{\xi r} & D_{\xi\theta}\\
D_{rp} & D_{r\xi} & D_{rr} & D_{r\theta}\\
D_{\theta p} & D_{\theta\xi} & D_{\theta r} & D_{\theta\theta}\end{array}\right)\]
 and \[
\overline{\overline{F}}=\left(\begin{array}{l}
F_{p}\\
F_{\xi}\\
F_{r}\\
F_{\theta}\end{array}\right)\]
 where each element is function of the momentum $p$ $\in\left[0,+\infty\right]$
and the cosine of the pitch-angle $\xi\in\left[-1,+1\right]$ in the
momentum space, but also depends of the radial position $r\in\left[0,a\right]$
and the polo\"{\i}dal angle $\theta\in\left[0,2\pi\right]$ in the
configuration space. All cross-coupling terms are neglected, hence
$D_{pr}=D_{rp}=0,$ $D_{\xi r}=D_{r\xi}=0,$ $D_{p\theta}=D_{\theta p}=0$
et $D_{\xi\theta}=D_{\theta\xi}=0.$ In these conditions, dynamics
in momentum and configuration spaces are decoupled, \[
\frac{\partial f}{\partial t}+\nabla_{p}\cdot\mathbf{S}_{\mathbf{p}}+\nabla_{r}\cdot\mathbf{S}_{\mathbf{r}}=0\]
 where

\[
\left\{ \begin{array}{l}
\mathbf{S}_{\mathbf{p}}=-\overline{\overline{D}}_{p}\nabla_{p}f_{0}+\overline{\overline{F}}_{p}f_{0}\\
\mathbf{S}_{\mathbf{r}}=-\overline{\overline{D}}_{r}\nabla_{r}f_{0}+\overline{\overline{F}}_{r}f_{0}\end{array}\right.\]
 the flux components being defined by the relations,

\[
\left\{ \begin{array}{l}
\mathbf{S}_{\mathbf{p}}=\left(\begin{array}{l}
S_{p}\\
S_{\xi}\end{array}\right)\\
\mathbf{S}_{\mathbf{r}}=\left(\begin{array}{l}
S_{r}\\
S_{\theta}\end{array}\right)\end{array}\right.\]


With this simplifying assumption, the diffusion tensor becomes \[
\overline{\overline{D}}=\left(\begin{array}{ll}
\overline{\overline{D}}_{p} & \overline{\overline{0}}\\
\overline{\overline{0}} & \overline{\overline{D}}_{r}\end{array}\right)=\left(\begin{array}{llll}
D_{pp} & D_{p\xi} & 0 & 0\\
D_{\xi p} & D_{\xi\xi} & 0 & 0\\
0 & 0 & D_{rr} & D_{r\theta}\\
0 & 0 & D_{\theta r} & D_{\theta\theta}\end{array}\right)\]
 while the convective term remains unchanged \[
\overline{\overline{F}}=\left(\begin{array}{l}
\overline{\overline{F}}_{p}\\
\overline{\overline{F}}_{r}\end{array}\right)=\left(\begin{array}{l}
F_{p}\\
F_{\xi}\\
F_{r}\\
F_{\theta}\end{array}\right)\]


Since the collision operator has a spherical symmetry, it is adequate
to choose the spherical coorodinate system for describing the electron
dynamic, since the numerical scheme is naturally well-conditionned
(dominant diagonal matrix). Considering this metric form,

\[
\nabla_{\mathbf{p}}\cdot\mathbf{S}_{\mathbf{p}}=\frac{1}{p^{2}}\frac{\partial\left(p^{2}S_{p}\right)}{\partial p}-\frac{1}{p}\frac{\partial\left(\sqrt{1-\xi^{2}}S_{\xi}\right)}{\partial\xi}\]
 where

\[
\nabla_{\mathbf{p}}f_{0}=\left(\begin{array}{l}
\frac{\partial f_{0}}{\partial p}\\
-\frac{\sqrt{1-\xi^{2}}}{p}\frac{\partial f_{0}}{\partial\xi}\end{array}\right)\]


Conversely, for the dynamic in the configuration space, the cylindrical
symmetry is more appropriate, and \[
\nabla_{\mathbf{r}}\cdot\mathbf{S}_{\mathbf{r}}=\frac{1}{r}\frac{\partial\left(rS_{r}\right)}{\partial r}+\frac{1}{r}\frac{\partial S_{\theta}}{\partial\theta}\]
 where \[
\nabla_{\mathbf{r}}f_{0}=\left(\begin{array}{l}
\frac{\partial f_{0}}{\partial r}\\
\frac{1}{r}\frac{\partial f_{0}}{\partial\theta}\end{array}\right)\]


The partial differential equation that must be solved is therefore,
\[
\frac{\partial f_{0}}{\partial t}+\frac{1}{p^{2}}\frac{\partial\left(p^{2}S_{p}\right)}{\partial p}-\frac{1}{p}\frac{\partial\left(\sqrt{1-\xi^{2}}S_{\xi}\right)}{\partial\xi}+\frac{1}{r}\frac{\partial\left(rS_{r}\right)}{\partial r}+\frac{1}{r}\frac{\partial S_{\theta}}{\partial\theta}=0\]
 where, in order to avoid spurious singularities at $p=r=0$, it is
useful to multiply both sides by $p^{2}$ and $r$,\[
\frac{\partial\left(rp^{2}f_{0}\right)}{\partial t}+r\frac{\partial\left(p^{2}S_{p}\right)}{\partial p}-rp\frac{\partial\left(\sqrt{1-\xi^{2}}S_{\xi}\right)}{\partial\xi}+p^{2}\frac{\partial\left(rS_{r}\right)}{\partial r}+p^{2}\frac{\partial S_{\theta}}{\partial\theta}=0\]
 The exact value of $f_{0}$ at $p=r=0$ is therefore never known
exactly, but just asymptotically approached, depending upon the size
of the mesh grid.


